\begin{figure}[htbp]
  \centering
  \begin{subfigure}[b]{0.31\linewidth}
    \centering
    \includegraphics[width=\linewidth]{benzene_HOMO.png}
    \caption{HOMO ($-6.97$~eV)}
  \end{subfigure}\hfill
  \begin{subfigure}[b]{0.31\linewidth}
    \centering
    \includegraphics[width=\linewidth]{benzene_HOMOm1.png}
    \caption{HOMO$-1$ ($-6.97$~eV)}
  \end{subfigure}\hfill
  \begin{subfigure}[b]{0.31\linewidth}
    \centering
    \includegraphics[width=\linewidth]{benzene_LUMO.png}
    \caption{LUMO ($-0.25$~eV)}
  \end{subfigure}
  \caption{Frontier orbitals of benzene (B3LYP/def2-SVP; isosurfaces at $\pm 0.03$~a.u.;
  blue and red denote opposite phases). (a) and (b) form the degenerate $e_{1g}$ pair, each with one
  nodal plane perpendicular to the ring; (c) belongs to the degenerate $e_{2u}$ pair, with two nodal
  planes.}
  \label{fig:benzene}
\end{figure}
